\edef\refsec{\refmodelsec}
\edef\refeq{\refmodeleq}

\section{Informal Savings Model}\label{\refsec}
This model is a variation of the informal analytical model in one respect: the informal households $j=0$ save, through an informal vehicle that earns $R_t^0 = \chi_t^RR_t$ and is not part of the formal capital stock, so informal savings do not enter the aggregate $s_t$. Every core object of the economic equilibrium sums over the formal types only, and the economic equilibrium is that of the analytical model except for the informal households' own choices.

\smalltitle{What changes in the political problem.}
The formal savings shares $s_{t-1,i}/s_{t-1}$ that the retirees carry into $t$ are functions of the tax $\tau_t$ they anticipated, so the policy maker's problem can substitute them out and they need not be states. The informal savings ratio $s_{t-1,0}/s_{t-1}$ admits no such reduction: it also depends on $\tau_{t-1}$, through formal savings $s_{t-1}$, which the policy maker at $t$ takes as given. It is therefore an endogenous state of the political problem, and the continuation policy is a function of it.


\subsection{The model}\label{\refsec:model}
We consider an economy populated by overlapping generations of workers and retirees. Within each generational cohort, we consider $J+1$ different household types indexed $j= 0,...,J$. We refer to type $j=0$ households as "informal"; they work in informal markets and save through an informal technology that is not part of the formal capital stock. Households of type $i>0$ are "formal" households; they supply $h_{t,i}\eta_{t,i}$ units of labour when young with $\eta_{t,i}$ denoting labour productivity. We let $\gamma_{t,j}$ denote sizes of working households with $\sum_i \gamma_{t,i} = 1$.

As young, households split a unit of time between labour and leisure. Households survive and retire with probability $p_{t,j}$; all formal households survive with the same probability $p_{t,i}= p_t$. Savings lost to mortality are insured within type, for informal as well as formal households.

\smalltitle{Technology.}
\begin{subequations}\label{\refeq:aggregateStates}
In the formal sector, a continuum of firms transforms aggregate capital and labour into output. We define aggregate labour and capital from
\begin{align}
    H_t =& n_th_t, & h_t&\equiv \sum_{i}\gamma_{t,i}\eta_{t,i}h_{t,i} \\
    K_t =& n_{t-1} s_{t-1}, & s_{t-1}&\equiv \sum_i\gamma_{t-1,i}s_{t-1,i}.
\end{align}
\end{subequations}
Both sums run over the formal types $i>0$ only: informal labour is supplied outside the formal sector, and informal savings $s_{t,0}$ are held in the informal savings vehicle and do \emph{not} add to the formal capital stock. The representative firm produces output with a Cobb-Douglas technology, with total factor productivity normalised to one,
\begin{align*}
    Y_t = K_t^{\alpha}H_t^{1-\alpha}.
\end{align*}
In equilibrium this implies the following factor prices:
\begin{align}\label{\refeq:factorPrices}
    R_t = \alpha \left(\dfrac{s_{t-1}}{\nu_t h_t}\right)^{\alpha-1}, \qquad w_t = (1-\alpha) \left(\dfrac{s_{t-1}}{\nu_t h_t}\right)^{\alpha}.
\end{align}


\bigskip
Informal households employ a linear production technology with return $w_t^0$ per efficiency unit of labour and $R_t^0$ per unit of informal savings. For simplicity, we model both as proportional to the corresponding formal factor price absent taxes writing
\begin{align}\label{\refeq:factorPrices0}
    w_t^0= \left(\frac{1-\alpha}{\Gamma_{h,t}^{\alpha}}\right)^{\frac{1}{1+\alpha\xi}} \left(\frac{s_{t-1}}{\nu_t}\right)^{\frac{\alpha}{1+\alpha\xi}}, \qquad R_t^0= \alpha \left(\frac{s_{t-1}}{\nu_th_t}\right)^{\alpha-1} \chi_t^R,
\end{align}
where $\chi_t^R$ scales the return relative to the formal capital market. We note that a similar scaling is not required for $w_t^0$ as the informal income already includes the type-specific productivity $\eta_{t,0}$.


\smalltitle{Preferences and Household Choices.}
Workers and retirees in period $t$ value consumption, $c_{1,t}^j$ and $c_{2,t}^j$ respectively. Workers discount the future at factor $\beta_{t,j}>0$, which includes the survival probability $p_{t,j}$.

For the young households, the expected utilities are given by
\begin{align}\label{\refeq:util}
    &\max\mbox{ } U_{t,j}\left(c_{1,t}^j, h_{t,j}, c_{2,t+1}^j\right) = u_{t,j}\left(c_{1,t}^j, h_{t,j}\right) +\beta_{t,j}u_{t,j}\left(c_{2,t+1}^j,0\right)
\end{align}
where we use CRRA-GHH preferences
\begin{align*}
    u_{t,j}(c,h) = \left\lbrace\begin{array}{ll} \dfrac{1}{1-1/\rho}\left(\tilde{c}_{t}^j\right)^{1-1/\rho}, & \text{for }\rho \neq 1 \\
    \ln\left(\tilde{c}_{t}^j\right) , & \text{for }\rho = 1
    \end{array}\right. , \qquad \tilde{c}_t^j \equiv c-\dfrac{\xi}{1+\xi}X_{t,j}\left(h\right)^{\frac{1+\xi}{\xi}}
\end{align*}
where $\rho>0$ is the intertemporal elasticity of substitution, $\xi$ is the Frisch elasticity, and $X_{t,j}$ is a parameter that measures the disutility of labour.

\begin{subequations}\label{\refeq:formalBudget}
The formal household receives wages, pays social security taxes $\tau_t$, saves $s_{t,i}$ and consumes when young, and consumes its savings and pension benefits $b_{t+1}^i$ when old:
\begin{align}
    c_{1,t}^i &= w_t\eta_{t,i} h_{t,i}(1-\tau_t)-s_{t,i} \\
    c_{2,t+1}^i &= s_{t,i}R_{t+1}/p_t+ b_{t+1}^i(h_{t,i}).
\end{align}
\end{subequations}

\begin{subequations}\label{\refeq:informalBudget}
Informal households earn informal wages when young, of which they consume part and save part in the informal vehicle; when old they consume their informal savings including pension benefits $b_{t+1}^0$:
\begin{align}
    c_{1,t}^0 &= w_t^0\eta_{t,0}h_{t,0} - s_{t,0} \\
    c_{2,t+1}^0 &= s_{t,0}R^0_{t+1}/p_{t,0}+b_{t+1}^0,
\end{align}
As for formal households, savings lost to mortality are insured within the type, which is the source of the $1/p_{t,0}$ term.\footnote{We do not formally restrict $s_{t,0}$ to be non-negative. In every calibration, however, $s_{t,0}/s_t$ lies inside the unit interval.}
\end{subequations}



\smalltitle{Social Security System.}
The government runs a pay-as-you-go pension system with a balanced budget. Every period, the government taxes labour income in the formal sector and transfers the sum directly to current retirees; informal savings are not taxed. Systems differ along two dimensions: whether retirees need to qualify for full benefits, and whether benefits are tied to earnings or flat-rate. Policy is set by a government that acts in the interest of voters, but lacks commitment.

Benefits are given by
\begin{subequations} \label{\refeq:governmentBudget}
\begin{align}
    b_t^i =& \left[\theta_th_{t-1,i} \eta_{t-1,i}+(1-\theta_t)h_{t-1}\right] \overline{b}_t \\
    b_t^0 =& \epsilon_t h_{t-1} \overline{b}_t
\end{align}
where $\epsilon_t$ measures universal pensions and $\theta_t$ measures contributive incentives.



We note that the shares $\sum_j\gamma_{t,j} = 1+\gamma_{t,0}$ to have a unit mass of formal households (convenient to have for definition of aggregate states). Thus, the average contribution of young households at time $t$ is
\begin{align*}
    \text{Avg. contribution from formal} = w_t\tau_t\sum_{i>0}\gamma_{t,i}\eta_{t,i}h_{t,i} = w_t\tau_th_t.
\end{align*}
Letting $n_t$ denote total number of young households, the number of formal households are $n_t/(1+\gamma_{t,0})$. Similarly, without mortality, the share of retired households of type $j$ would be $\gamma_{t-1,j}/(1+\gamma_{t-1,0})$. Taking mortality into consideration, the balanced budget requires that
\begin{align*}
    \dfrac{n_t}{1+\gamma_{t,0}} w_t\tau_th_t = \dfrac{n_{t-1}}{1+\gamma_{t-1,0}}\sum_j \gamma_{t-1,j}p_{t-1,j} b_t^j
\end{align*}
The left hand side measures contributions, the right hand side measures beneficiaries. Using the structure of the pension design, this identifies the level in pension rates $\overline{b}_t$ as:
\begin{align}\label{\refeq:governmentBudget:bbar}
    \overline{b}_t & \equiv \dfrac{\nu_t w_th_t\tau_t}{h_{t-1}\kappa_{t-1}}, \qquad\kappa_t \equiv \frac{1+\gamma_{t,0}}{1+\gamma_{t+1,0}}(\epsilon_{t+1}\gamma_{t,0}p_{t,0}+p_t)
\end{align}
\end{subequations}


\bigskip
\begin{subequations}\label{\refeq:formalOpt}
With CRRA-GHH preferences, households maximize \eqref{\refeq:util} subject to budgets \eqref{\refeq:formalBudget} and \eqref{\refeq:informalBudget}. The labour supply and savings of formal households are then given by (given factor prices and policies):
\begin{align}
    h_{t,i} &= \left[\frac{\eta_{t,i}}{X_{t,i}}\left(w_t(1-\tau_t)+\frac{\theta_{t+1}\overline{b}_{t+1}}{R_{t+1}/p_t}\right)\right]^{\xi} \\
    s_{t,i} &= \dfrac{B_{t+1}^i}{1+B_{t+1}^i}\left[w_t\left(1-\tau_t\right)h_{t,i}\eta_{t,i}-\dfrac{X_{t,i}\xi}{1+\xi}\left(h_{t,i}\right)^{\frac{1+\xi}{\xi}}\right]-\dfrac{b_{t+1}^i/(R_{t+1}/p_t)}{1+B_{t+1}^i}
\end{align}
\end{subequations}
where we have used the shorthand $B_{t+1}^i \equiv \left(\beta_{t,i}\right)^{\rho}(R_{t+1}/p_t)^{\rho-1}$. An hour worked at $t$ raises the household's own benefit at $t+1$ by $\theta_{t+1}\eta_{t,i}\overline{b}_{t+1}$, which the Euler equation discounts at $R_{t+1}/p_t$.

For the informal households, similar conditions apply with a few adjustments for relevant factor prices and the incentive structure of universal pensions:
\begin{subequations}\label{\refeq:informalOpt}
\begin{align}
    h_{t,0} &= \left(\frac{\eta_{t,0}}{X_{t,0}}w_t^0\right)^{\xi},\\
    s_{t,0} &=  \dfrac{B_{t+1}^0}{1+B_{t+1}^0}\left[w_t^0h_{t,0}\eta_{t,0}-\dfrac{X_{t,0}\xi}{1+\xi}\left(h_{t,0}\right)^{\frac{1+\xi}{\xi}}\right]-\dfrac{b_{t+1}^0/(R_{t+1}^0/p_{t,0})}{1+B_{t+1}^0},
\end{align}
\end{subequations}
where $B_{t+1}^0 \equiv \left(\beta_{t,0}\right)^{\rho} \left(R_{t+1}^0/p_{t,0}\right)^{\rho-1}$ is the generalized discount function for the informal households.



\subsection{The Competitive Equilibrium}\label{\refsec:EE}
In the competitive equilibrium firms and households maximize, the government runs a balanced budget, and the markets clear. We can collect these in definition of aggregate states \eqref{\refeq:aggregateStates}, factor prices \eqref{\refeq:factorPrices}, household budgets \eqref{\refeq:formalBudget} and \eqref{\refeq:informalBudget}, government budgets \eqref{\refeq:governmentBudget}, and household choices \eqref{\refeq:formalOpt} and \eqref{\refeq:informalOpt}. The following derives closed form representations of the core model variables \emph{given} the generalized discount factor $B_{t+1}^j$. 


\bigskip
\begin{subequations}\label{\refeq:EE}
Using equilibrium conditions for wages, interest rates, and pension benefits, the labour supply for type $i$ can in equilibrium be written as
\begin{align*}
    h_{t,i} = \left(\dfrac{\eta_{t,i}}{X_{t,i}}\right)^{\xi}\left(\dfrac{1}{h_t}\right)^{\xi}\left((1-\alpha)(1-\tau_t)\left(\dfrac{s_{t-1}}{\nu_t}\right)^{\alpha}h_t^{1-\alpha}+\dfrac{1-\alpha}{\alpha}\dfrac{p_t \theta_{t+1}}{\kappa_t} s_t\tau_{t+1}\right)^{\xi}.
\end{align*}
The only type-specific term is $(\eta_{t,i}/X_{t,i})^{\xi}$; this shows that the ratio $h_{t,i}/h_t$ is only a function of parameters, namely that
\begin{align}\label{\refeq:EE:hi}
    \dfrac{h_{t,i}}{h_t} = \dfrac{(\eta_{t,i}/X_{t,i})^{\xi}}{\Gamma_{h,t}}, \qquad \Gamma_{h,t} = \sum_i\frac{\gamma_{t,i}\eta_{t,i}^{1+\xi}}{X_{t,i}^{\xi}}.
\end{align}


Next, we can use the wage rate and individual labour supply functions to write
\begin{align*}
    w_t(1-\tau_t)h_{t,i}\eta_{t,i}-\dfrac{X_{t,i}\xi}{1+\xi}\left(h_{t,i}\right)^{\frac{1+\xi}{\xi}} = \dfrac{\eta_{t,i}^{1+\xi}}{X_{t,i}^{\xi}}\frac{1}{\Gamma_{h,t}}\left[(1-\alpha)(1-\tau_t)h_t^{1-\alpha}\left(\dfrac{s_{t-1}}{\nu_t}\right)^{\alpha}-\dfrac{\xi}{1+\xi}\frac{h_t^{\frac{1+\xi}{\xi}}}{\Gamma_{h,t}^{1/\xi}}\right]
\end{align*}
Using the expression for $h_{t,i}$ and $h_{t,i}/h_t$, we can rewrite this as:
\begin{align*}
    (1-\alpha)(1-\tau_t)h_t^{1-\alpha}\left(\dfrac{s_{t-1}}{\nu_t}\right)^{\alpha} = \frac{h_t^{\frac{1+\xi}{\xi}}}{\Gamma_{h,t}^{1/\xi}} - \dfrac{1-\alpha}{\alpha}\dfrac{p_t \theta_{t+1}}{\kappa_t} s_t\tau_{t+1},
\end{align*}
From this and discounted pension benefits in equilibrium, we then finally have that
\begin{align}\label{\refeq:EE:si}
    s_{t,i} =& \dfrac{1}{1+B_{t+1}^i}\left[B_{t+1}^i\frac{\eta_{t,i}^{1+\xi}}{X_{t,i}^{\xi}}\left(\frac{h_t}{\Gamma_{h,t}}\right)^{\frac{1+\xi}{\xi}}\frac{1}{1+\xi}-\frac{1-\alpha}{\alpha}\frac{p_t\left(1-\theta_{t+1}\right)}{\kappa_t}s_t\tau_{t+1}\right]\\
    -&\frac{\eta_{t,i}^{1+\xi}}{X_{t,i}^{\xi}\Gamma_{h,t}}\frac{1-\alpha}{\alpha}\frac{p_t\theta_{t+1}}{\kappa_t}s_t\tau_{t+1} \notag
\end{align}
Aggregating over the different types, the savings state can then be defined as
\begin{align}\label{\refeq:EE:s(h)}
    s_t = \left(\frac{h_t}{\Gamma_{h,t}}\right)^{\frac{1+\xi}{\xi}}\underbrace{\frac{1}{1+\xi}\frac{\sum_i \frac{\gamma_{t,i}\eta_{t,i}^{1+\xi}}{X_{t,i}^{\xi}}\frac{B_{t+1}^i}{1+B_{t+1}^i}}{1+\frac{1-\alpha}{\alpha}\frac{p_t\tau_{t+1}}{\kappa_t}\left(\theta_{t+1}+(1-\theta_{t+1})\sum_i\frac{\gamma_{t,i}}{1+B_{t+1}^i}\right)}}_{\equiv \Gamma_{s,t}}
\end{align}
Importantly, note that the savings ratios $s_{t,i}/s_t$ can be written as a function of parameters and \emph{future} policy $\tau_{t+1}$ (independent e.g. of $\tau_t$ and given the generalized discount function $B$). Formally we have:
\begin{align}\label{\refeq:EE:si_s}
    \frac{s_{t,i}}{s_t} = &\frac{\frac{B_{t+1}^i}{1+B_{t+1}^i}\frac{\eta_{t,i}^{1+\xi}}{X_{t,i}^{\xi}}}{\sum_{j>0} \frac{\gamma_{t,j}\eta_{t,j}^{1+\xi}}{X_{t,j}^{\xi}}\frac{B_{t+1}^j}{1+B_{t+1}^j}}\left[1+\frac{1-\alpha}{\alpha}\frac{p_t\tau_{t+1}}{\kappa_t}\left(\theta_{t+1}+(1-\theta_{t+1})\sum_{j>0}\frac{\gamma_{t,j}}{1+B_{t+1}^j}\right)\right] \\
    -&\frac{1}{1+B_{t+1}^i}\frac{1-\alpha}{\alpha}\frac{p_t(1-\theta_{t+1})}{\kappa_t}\tau_{t+1} - \frac{\eta_{t,i}^{1+\xi}/X_{t,i}^{\xi}}{\Gamma_{h,t}}\frac{1-\alpha}{\alpha}\frac{p_t\theta_{t+1}}{\kappa_t}\tau_{t+1}\notag
\end{align}
We note that when discount rates are identical across types, i.e. when $\beta_{t,i} = \beta_t$ and $B_{t+1}^i = B_{t+1}$, this simplifies to 
\begin{align*}
    \frac{s_{t,i}}{s_t} = \frac{\eta_{t,i}^{1+\xi}/X_{t,i}^{\xi}}{\Gamma_{h,t}}+\frac{1-\alpha}{\alpha}\frac{p_t\tau_{t+1}}{\kappa_t}\frac{1-\theta_{t+1}}{1+B_{t+1}}\left(\frac{\eta_{t,i}^{1+\xi}/X_{t,i}^{\xi}}{\Gamma_{h,t}}-1\right).
\end{align*}


Finally, substituting \eqref{\refeq:EE:s(h)} into $h_t = \sum_i\gamma_{t,i}\eta_{t,i}h_{t,i}$, with $h_{t,i}$ from the labour supply above, yields a closed form given the generalized discount functions:
\begin{align}\label{\refeq:EE:h}
    h_t = \underbrace{\left(\Gamma_{h,t}\right)^{\frac{1+\xi}{1+\alpha\xi}}\left(\frac{(1-\alpha)(1-\tau_t)}{\Gamma_{h,t}-\frac{1-\alpha}{\alpha}\frac{p_t\theta_{t+1}\tau_{t+1}}{\kappa_t}\Gamma_{s,t}}\right)^{\frac{\xi}{1+\alpha\xi}}}_{\equiv \Theta_{h,t}}\left(\frac{s_{t-1}}{\nu_t}\right)^{\frac{\alpha\xi}{1+\alpha\xi}}.
\end{align}
Naturally, we can use this in the expression for $s_t$ to write savings as
\begin{align}\label{\refeq:EE:s}
    s_t &= \underbrace{\left(\frac{\Theta_{h,t}}{\Gamma_{h,t}}\right)^{\frac{1+\xi}{\xi}}\Gamma_{s,t}}_{\equiv\Theta_{s,t}}\left(\frac{s_{t-1}}{\nu_t}\right)^{\frac{\alpha(1+\xi)}{1+\alpha\xi}}.
\end{align}


Given discount functions, $B^i$, equations \eqref{\refeq:EE} identify each aggregate state as $x_t = \Theta_{x,t}(s_{t-1}/\nu_t)^{\sigma_x}$, a coefficient function $\Theta_{x,t}$ of parameters and policy times a power of the lagged state, with $\sigma_s = \alpha(1+\xi)/(1+\alpha\xi)$ and $\sigma_h = \alpha\xi/(1+\alpha\xi)$ independent of policy.


\bigskip
Given the structure of the aggregate states in economic equilibrium, it is straightforward to reduce the informal savings ratio $s_{t,0}/s_t$ and write it as
\begin{align}\label{\refeq:EE:s0_s}
    \frac{s_{t,0}}{s_t} =& \frac{B_{t+1}^0}{1+B_{t+1}^0}\frac{\eta_{t,0}^{1+\xi}}{X_{t,0}^{\xi}(1+\xi)}\left(\frac{1-\alpha}{\Gamma_{h,t}^{\alpha}}\right)^{\frac{1+\xi}{1+\alpha\xi}}
    \frac{1}{\Theta_{s,t}}-\frac{1}{1+B_{t+1}^0}\frac{1-\alpha}{\chi_{t+1}^R\alpha}\frac{p_{t,0}\epsilon_{t+1}\tau_{t+1}}{\kappa_t}.
\end{align}
\end{subequations}


\smalltitle{Closed form consumption functions.}
We note that if we know $s_t$, $h_t$, $B^i$, and $s_{t-1}$, it is straightforward to identify all other variables directly: Individual labour supply and savings decisions from \eqref{\refeq:EE:hi} and \eqref{\refeq:EE:si}, factor prices from \eqref{\refeq:factorPrices}, pension benefits from \eqref{\refeq:governmentBudget}, and finally consumption simply from budgets \eqref{\refeq:formalBudget} and \eqref{\refeq:informalBudget}. However, in particular in the analytical log case, it can be useful to derive \emph{equilibrium} consumption functions as well. All consumption functions take the same form $x_t = \Theta_{x,t}(s_{t-1}/\nu_t)^{\sigma_x}$, with
\begin{align}\label{\refeq:EE:sigma_ci}
    \sigma_{c_1,t}^i &= \sigma_{c_2,t}^i = \sigma_{\tilde{c}_1,t}^i = \dfrac{\alpha(1+\xi)}{1+\alpha\xi}, 
    \qquad \sigma_{c_2,t+1}^i = \left(\dfrac{\alpha(1+\xi)}{1+\alpha\xi}\right)^2.
\end{align}
This is essentially what ensures that the political objective function is log-separable in the state $s_{t-1}/\nu_t$. 

\begin{subequations}\label{\refeq:EE:ci}
The coefficient functions (that also generally depends on policy) are derived using the aggregate states in the functions above. For completeness, we have the following:
\begin{align}
    c_{1,t}^i &= \frac{\eta_{t,i}^{1+\xi}}{X_{t,i}^{\xi}}\left(\frac{h_t}{\Gamma_{h,t}}\right)^{\frac{1+\xi}{\xi}}\left(1-\frac{B_{t+1}^i}{(1+B_{t+1}^i)(1+\xi)}\right)+\frac{s_t}{1+B_{t+1}^i}\frac{1-\alpha}{\alpha}\frac{p_t \tau_{t+1}(1-\theta_{t+1})}{\kappa_t} \\
    c_{2,t}^i &= \alpha \frac{\nu_t}{p_{t-1}}h_t^{1-\alpha}\left(\frac{s_{t-1}}{\nu_t}\right)^{\alpha}\left(\frac{s_{t-1,i}}{s_{t-1}}+\frac{1-\alpha}{\alpha}\frac{p_{t-1}\tau_t}{\kappa_{t-1}}\left(1+\theta_t\left[\frac{\eta_{t-1,i}^{1+\xi}}{X_{t-1,i}^{\xi}}\frac{1}{\Gamma_{h,t-1}}-1\right]\right)\right) \\
    \tilde{c}_{1,t}^i &= \frac{\eta_{t,i}^{1+\xi}}{X_{t,i}^{\xi}}\left(\frac{h_t}{\Gamma_{h,t}}\right)^{\frac{1+\xi}{\xi}}\frac{1}{(1+\xi)(1+B_{t+1}^i)}+\frac{s_t}{1+B_{t+1}^i}\frac{1-\alpha}{\alpha}\frac{p_t \tau_{t+1}(1-\theta_{t+1})}{\kappa_t} \notag \\ 
    &= \left(\frac{h_t}{\Gamma_{h,t}}\right)^{\frac{1+\xi}{\xi}}\frac{1}{1+B_{t+1}^i}\left(\frac{\eta_{t,i}^{1+\xi}}{X_{t,i}^{\xi}}\frac{1}{1+\xi}+\Gamma_{s,t}\frac{1-\alpha}{\alpha}\frac{p_t \tau_{t+1}(1-\theta_{t+1})}{\kappa_t}\right)
\end{align}
For the formal households that do not supply labour in retirement, we simply have $\tilde{c}_{2,t}^i = c_{2,t}^i$. \end{subequations}


For the informal households that save through the informal vehicle, relevant equilibrium consumption functions are given by
\begin{subequations}\label{\refeq:EE:c0}
\begin{align}
    c_{1,t}^0 =& \dfrac{\eta_{t,0}^{1+\xi}}{X_{t,0}^{\xi}}(w_t^0)^{1+\xi}\left(1-\frac{B_{t+1}^0}{(1+B_{t+1}^0)(1+\xi)}\right)+\frac{1}{1+B_{t+1}^0}\frac{1-\alpha}{\chi_{t+1}^R \alpha}\frac{p_{t,0}\epsilon_{t+1}\tau_{t+1}}{\kappa_t} s_t \\ 
    \tilde{c}_{1,t}^0 =& \dfrac{\eta_{t,0}^{1+\xi}}{X_{t,0}^{\xi}}\frac{(w_t^0)^{1+\xi}}{(1+B_{t+1}^0)(1+\xi)}+\frac{1}{1+B_{t+1}^0}\frac{1-\alpha}{\chi_{t+1}^R \alpha}\frac{p_{t,0}\epsilon_{t+1}\tau_{t+1}}{\kappa_t} s_t
\end{align}
Consumption when old is then given by
\begin{align}
    c_{2,t}^0 =& \chi_{t}^R\alpha \frac{\nu_t}{p_{t-1,0}}\left(\frac{s_{t-1}}{\nu_t}\right)^{\alpha}h_t^{1-\alpha}\left(\frac{s_{t-1,0}}{s_{t-1}}+\frac{1-\alpha}{\chi_{t}^R\alpha} \frac{p_{t-1,0}\epsilon_t\tau_t}{\kappa_{t-1}}\right).
\end{align}
\end{subequations}



\subsection{The politico-economic equilibrium with LOG preferences}\label{\refsec:PEELOG}
\begin{subequations}\label{\refeq:PEELOG}
Let $\upsilon_{1,t}^j$ and $\upsilon_{2,t}^j$ denote the indirect utilities of different households (types and generations). Policy is set without commitment, with elections each $t$ aggregating preferences using probabilistic voting that maximizes a weighted sum of indirect utilities. Let $\tau^{t+1}(z_t)$ denote a continuation policy with $z_t$ being a vector of states. The political problem can then be written as
\begin{align*}
    &\max_{\tau_t}\mbox{ }\mathcal{W}_t(z_{t-1}) \equiv \omega\sum_j\left(\gamma_{t-1,j} p_{t-1,j} \mu_{t-1,j} \upsilon_{2,t}^j\right)+\nu_t\sum_j\left(\gamma_{t,j}\mu_{t,j} \upsilon_{1,t}^j\right) \\
    &\text{s.t. competitive equilibrium and }\tau_{t+1} = \tau^{t+1}(z_t),
\end{align*}
where $\tau^{t+1}(z_t)$ is identified by this process at $t+1$, $\omega$ is a weight reflecting relative political power of the older generation, $\mu_{t,j}$ reflects relative voting shares for different income groups (the old generation at $t$ votes with the voting shares $\mu_{t-1,j}$ they had when young, i.e. shares are attached to a generation, not a period). In general, the first order condition for $\tau_t$ is given by:
\begin{align*}
    \omega\sum_{j}\left(\gamma_{t-1,j}p_{t-1,j}\mu_{t-1,j} \frac{d\upsilon_{2,t}^j}{d\tau_t}\right)+
    \nu_t\sum_j\left(\gamma_{t,j}\mu_{t,j} \frac{d\upsilon_{1,t}^j}{d\tau_t}\right) = 0,
\end{align*}
where the total derivatives (with "hard d") includes the partial effects of changing policy directly (through economic equilibrium functions), but also the indirect effects working through the continuation policy. 

With $\rho = 1$, the discount function simplifies to $B_{t+1}^i = \beta_{t,i}$ and the continuation policy depends only on the endogenous state $s_{t,0}/s_t$.
For the consumption smoothers, using the Euler condition, we have $\upsilon_{1,t}^i = (1+\beta_{t,i})\ln\left(\tilde{c}_{1,t}^i\right)+\beta_{t,i} \ln\left(R_{t+1}\right)+\beta_{t,i} \ln(\beta_{t,i}/p_{t,i})$. The marginal effect of taxes can then be written as
\begin{align}
    \frac{d\upsilon_{1,t}^i}{d\tau_t} = (1+\beta_{t,i})\frac{d\ln\left(\tilde{c}_{1,t}^i\right)}{d\tau_t} + \beta_{t,i} \frac{d\ln\left(R_{t+1}\right)}{d\tau_t}.
\end{align}
For the informal young, we can similarly use an Euler condition to simplify this as $\upsilon_{1,t}^0 = \left(1+\beta_{t,0}\right)\ln\left(\tilde{c}_{1,t}^0\right)+\beta_{t,0} \ln\left(R_{t+1}\right)+\beta_{t,0}\ln\left(\beta_{t,0}\chi_{t+1}^R/p_{t,0}\right)$, thus we have
\begin{align}
    \frac{d\upsilon_{1,t}^0}{d\tau_t} = (1+\beta_{t,0})\frac{d\ln\left(\tilde{c}_{1,t}^0\right)}{d\tau_t} + \beta_{t,0} \frac{d\ln\left(R_{t+1}\right)}{d\tau_t}.
\end{align}



For the older cohorts, the effect works through $\tilde{c}_{2,t}^j$ and, again being log-separable in $s_{t-1}$, this simplifies to
\begin{align}
    \frac{d\upsilon_{2,t}^i}{d\tau_t} =& (1-\alpha)\frac{d\ln\left(h_t\right)}{d \tau_t}+\frac{\frac{1-\alpha}{\alpha}\frac{p_{t-1}}{\kappa_{t-1}}\left(1+\theta_t\left[\frac{\eta_{t-1,i}^{1+\xi}}{X_{t-1,i}^{\xi}}\frac{1}{\Gamma_{h,t-1}}-1\right]\right)}{\frac{s_{t-1,i}}{s_{t-1}}+\frac{1-\alpha}{\alpha}\frac{p_{t-1}\tau_t}{\kappa_{t-1}}\left(1+\theta_t\left[\frac{\eta_{t-1,i}^{1+\xi}}{X_{t-1,i}^{\xi}}\frac{1}{\Gamma_{h,t-1}}-1\right]\right)} \\
    \frac{d\upsilon_{2,t}^0}{d\tau_t} =& (1-\alpha)\frac{d\ln\left(h_t\right)}{d \tau_t}+\frac{\frac{1-\alpha}{\chi_t^R\alpha}\frac{p_{t-1,0}\epsilon_t}{\kappa_{t-1}}}{\frac{s_{t-1,0}}{s_{t-1}}+\frac{1-\alpha}{\chi_t^R\alpha}\frac{p_{t-1,0}\epsilon_t\tau_t}{\kappa_{t-1}}}.
\end{align}
As we show in \cref{\refnumsec}, we will rely on numerical methods that approximate these derivatives in combination with grid searches.
\end{subequations}





\subsection{The politico-economic equilibrium with CRRA preferences}\label{\refsec:PEE}
With CRRA preferences, the general first order conditions largely remain the same, but with consumption levels now playing a role and total derivatives including effects through both endogenous state variables now. 

\begin{subequations}\label{\refeq:PEE}
Specifically, we now have for all household types (including $j=0$):
\begin{align}
    \upsilon_{1,t}^j = (1+B_{t+1}^j)\frac{\left(\tilde{c}_{1,t}^j\right)^{1-1/\rho}}{1-1/\rho}
\end{align}


Because the policy maker has to take past choices as given (such as savings and relative savings rates), we have to rely on first order conditions to efficiently identify their choices (to avoid searching over a full grid of states $s_{t-1,i}/s_{t-1}$). Whereas we can numerically approximate derivatives of $\upsilon_{1,t}^i$  and $\upsilon_{1,t}^0$ using simple grids, here we have to be a bit more careful. In particular, we have to use:
\begin{align}
    \dfrac{d \upsilon_{2,t}^i}{d \tau_t} =& \left(c_{2,t}^i\right)^{1-1/\rho}\dfrac{d\ln\left(c_{2,t}^i\right)}{d\tau_t} \\
    \frac{d\ln\left(c_{2,t}^i\right)}{d\tau_t} =& (1-\alpha)\frac{d \ln\left(h_t\right)}{d \tau_t}+\frac{\frac{1-\alpha}{\alpha}\frac{p_{t-1}}{\kappa_{t-1}}\left(1+\theta_t\left[\frac{\eta_{t-1,i}^{1+\xi}}{X_{t-1,i}^{\xi}}\frac{1}{\Gamma_{h,t-1}}-1\right]\right)}{\frac{s_{t-1,i}}{s_{t-1}}+\frac{1-\alpha}{\alpha}\frac{p_{t-1}\tau_t}{\kappa_{t-1}}\left(1+\theta_t\left[\frac{\eta_{t-1,i}^{1+\xi}}{X_{t-1,i}^{\xi}}\frac{1}{\Gamma_{h,t-1}}-1\right]\right)},
\end{align}
where the second line explicitly uses that $s_{t-1,i}/s_{t-1}$ is predetermined. However, as this ratio has a closed-form expression that relies on $\tau_t$, we do not have to evaluate on the entire grid, rather only on the selection that are actually consistent with economic equilibrium conditions. This is the exact trick that allows us to only use $s_{t-1}$ as a relevant state and not the entire distribution.


Finally, for the old informal households, we similarly have to rely on the first order condition:
\begin{align}
    \dfrac{d \upsilon_{2,t}^0}{d \tau_t} &= \left(c_{2,t}^0\right)^{1-1/\rho} \dfrac{d\ln\left(c_{2,t}^0\right)}{d\tau_t} \\
    \dfrac{d\ln\left(c_{2,t}^0\right)}{d\tau_t} &= (1-\alpha)\frac{d\ln\left(h_t\right)}{d \tau_t}+\frac{\frac{1-\alpha}{\chi_t^R\alpha}\frac{p_{t-1,0}\epsilon_t}{\kappa_{t-1}}}{\frac{s_{t-1,0}}{s_{t-1}}+\frac{1-\alpha}{\chi_t^R\alpha}\frac{p_{t-1,0}\epsilon_t\tau_t}{\kappa_{t-1}}}.
\end{align}
\end{subequations}
\subsection{The terminal period}\label{\refsec:finitehorizon}
\smalltitle{Economic equilibrium.}
In the finite horizon version of the model, the economic equilibrium simplifies to (with $s_t = s_{t,j} = 0$):
\begin{subequations}\label{\refeq:TerminalEE}
\begin{align}
    h_t &= \left(\Gamma_{h,t}\right)^{\frac{1}{1+\alpha\xi}}\left((1-\alpha)(1-\tau_t)\right)^{\frac{\xi}{1+\alpha\xi}}\left(\frac{s_{t-1}}{\nu_t}\right)^{\frac{\alpha\xi}{1+\alpha\xi}} \\
    h_{t,i} &= h_t\frac{(\eta_{t,i}/X_{t,i})^{\xi}}{\Gamma_{h,t}} \\
    c_{1,t}^i &= \dfrac{\eta_{t,i}^{1+\xi}}{X_{t,i}^{\xi}}\left(\frac{h_t}{\Gamma_{h,t}}\right)^{\frac{1+\xi}{\xi}} \\
    \tilde{c}_{1,t}^i &= \frac{c_{1,t}^i}{1+\xi} \\
    c_{1,t}^0 &= \frac{\eta_{t,0}^{1+\xi}}{X_{t,0}^{\xi}}\left(\frac{1-\alpha}{\Gamma_{h,t}^{\alpha}}\right)^{\frac{1+\xi}{1+\alpha\xi}}\left(\frac{s_{t-1}}{\nu_t}\right)^{\frac{\alpha(1+\xi)}{1+\alpha\xi}} \\
    \tilde{c}_{1,t}^0 &= \dfrac{c_{1,t}^0}{1+\xi}.
\end{align}
\end{subequations}

\smalltitle{Politico-economic equilibrium with log preferences.}
Without a continuation policy, the marginal indirect utilities of the terminal period are closed form:
\begin{subequations}\label{\refeq:terminalPEELOG}
\begin{align}
    \dfrac{d \upsilon_{1,t}^i}{d \tau_t} =& - \frac{1}{1-\tau_t}\frac{1+\xi}{1+\alpha\xi} \\
    \dfrac{d \upsilon_{1,t}^0}{d \tau_t} =& 0 \\
    \frac{d\upsilon_{2,t}^i}{d\tau_t} =& (1-\alpha)\frac{d\ln\left(h_t\right)}{d \tau_t}+\frac{\frac{1-\alpha}{\alpha}\frac{p_{t-1}}{\kappa_{t-1}}\left(1+\theta_t\left[\frac{\eta_{t-1,i}^{1+\xi}}{X_{t-1,i}^{\xi}}\frac{1}{\Gamma_{h,t-1}}-1\right]\right)}{\frac{s_{t-1,i}}{s_{t-1}}+\frac{1-\alpha}{\alpha}\frac{p_{t-1}\tau_t}{\kappa_{t-1}}\left(1+\theta_t\left[\frac{\eta_{t-1,i}^{1+\xi}}{X_{t-1,i}^{\xi}}\frac{1}{\Gamma_{h,t-1}}-1\right]\right)} \\
    \frac{d\upsilon_{2,t}^0}{d\tau_t} =& (1-\alpha)\frac{d\ln\left(h_t\right)}{d \tau_t}+\frac{\frac{1-\alpha}{\chi_t^R\alpha}\frac{p_{t-1,0}\epsilon_t}{\kappa_{t-1}}}{\frac{s_{t-1,0}}{s_{t-1}}+\frac{1-\alpha}{\chi_t^R\alpha}\frac{p_{t-1,0}\epsilon_t\tau_t}{\kappa_{t-1}}}.
\end{align}
where the $d \ln\left(h_t\right)/d \tau_t = -\frac{\xi}{1+\alpha\xi}\frac{1}{1-\tau_t}$ for this terminal state only.
\end{subequations}

\smalltitle{Politico-economic equilibrium with CRRA preferences.}
With $\rho\neq 1$ each term of the log condition is weighted by its consumption level to the power $1-1/\rho$, so the savings level $s_{t-1}$ becomes a relevant state:
\begin{subequations}\label{\refeq:terminalPEECRRA}
\begin{align}
    \dfrac{d \upsilon_{1,t}^i}{d \tau_t} =& - \frac{1}{1-\tau_t}\frac{1+\xi}{1+\alpha\xi}\left(\tilde{c}_{1,t}^i\right)^{1-1/\rho}  \\
    \dfrac{d \upsilon_{1,t}^0}{d \tau_t} =& 0 \\
    \frac{d\upsilon_{2,t}^i}{d\tau_t} =& \left(c_{2,t}^i\right)^{1-1/\rho}\left[(1-\alpha)\frac{d\ln\left(h_t\right)}{d \tau_t}+\frac{\frac{1-\alpha}{\alpha}\frac{p_{t-1}}{\kappa_{t-1}}\left(1+\theta_t\left[\frac{\eta_{t-1,i}^{1+\xi}}{X_{t-1,i}^{\xi}}\frac{1}{\Gamma_{h,t-1}}-1\right]\right)}{\frac{s_{t-1,i}}{s_{t-1}}+\frac{1-\alpha}{\alpha}\frac{p_{t-1}\tau_t}{\kappa_{t-1}}\left(1+\theta_t\left[\frac{\eta_{t-1,i}^{1+\xi}}{X_{t-1,i}^{\xi}}\frac{1}{\Gamma_{h,t-1}}-1\right]\right)} \right]\\
    \frac{d\upsilon_{2,t}^0}{d\tau_t} =& \left(c_{2,t}^0\right)^{1-1/\rho}\left[(1-\alpha)\frac{d\ln\left(h_t\right)}{d \tau_t}+\frac{\frac{1-\alpha}{\chi_t^R\alpha}\frac{p_{t-1,0}\epsilon_t}{\kappa_{t-1}}}{\frac{s_{t-1,0}}{s_{t-1}}+\frac{1-\alpha}{\chi_t^R\alpha}\frac{p_{t-1,0}\epsilon_t\tau_t}{\kappa_{t-1}}}\right],
\end{align}
where again $d \ln\left(h_t\right)/d \tau_t = -\frac{\xi}{1+\alpha\xi}\frac{1}{1-\tau_t}$ for this terminal state only.
\end{subequations}



\subsection{Calibration}\label{\refsec:calibration}
This model is only used to match the case of Argentina. The full procedure and a description of the data can be found in the main paper. The main application in the paper uses constant parameter values across $t$ for almost all parameters. Thus, when we drop the $t$-subscript in the following, we assume constant parameters.

\smalltitle{Data.}
We follow previous literature and let $t$ represent 30-year intervals (this helps with the assumption that capital fully depreciates between time steps). We find population data including projections to determine $\nu_t$ from 1950 to 2070, amend the data with $t = 1920, 2100$ by constant extrapolation, and extend the horizon by three further periods, to 2190, at the last value.

We use a household survey to split up the population into $J+1 = 5$ household types. This household survey provides shares $\gamma$, hours worked, and income. In addition, we have data on voting shares represented by $\mu_j$ in the political problem and survival rates $(p_j)$ (where $p_i = p$ is common across $i>0$ types).

We use the capital income share to identify $\alpha$, set the labour supply elasticity at a baseline value $\xi = 0.3$, and take the capital--output ratio to be 3.23, Argentina's ratio in the Penn World Table at the baseline year $t=2010$, where the remaining targets are measured too. Because a period is thirty years, the ratio the data report is thirty times the model's own $K_t/Y_t$, which is a stock over a single period's output; \cref{\refnumsec:calibration} sets out how that is imposed.

The tax target is derived from pension spending of $S = 7.1\%$ of GDP in the baseline year $t_0 = 2010$, with GDP the sum of formal and informal output. Spending is $\tau_{t_0}(1-\alpha)$ of formal output, and the informal income target below puts informal output at $a(1-\tau_{t_0})$ of formal output, with $a\equiv\gamma_0z_0^{\eta}(1-\alpha)$. The informal share of GDP is then $y^0 = a(1-\tau_{t_0})/(1+a(1-\tau_{t_0}))$, and $\tau_{t_0} = S/[(1-\alpha)(1-y^0)]$ solves in closed form,
\begin{align*}
    \tau_{t_0} = \frac{S(1+a)}{(1-\alpha)+Sa},
\end{align*}
which is $0.118$ in the Argentine data. On the solved path, spending over GDP equals $S$ wherever the informal targets are met.

We use OECD data to define ratio of replacement rates $RR$ at half mean wages and mean wages (represented by groups $j=1, j=3$ in our data). This will help us identify $\theta$ (we do not work with endogenous pension characteristics for Argentina).

The pre-reform universal pension is computed rather than calibrated. Before the reform, the basic pension paid 70\% of the minimum pension, read as the benefit of the least productive formal type $j=1$, and started five (men) or ten (women) years after the formal retirement age, so that
\begin{align*}
    \epsilon = 0.7\left(1-\theta+\theta\frac{\eta_1h_{t-1,1}}{h_{t-1}}\right)\frac{1}{2}\left(\beta^{5/30}\frac{\mathrm{LE_{men}}-5}{\mathrm{LE_{men}}}+\beta^{10/30}\frac{\mathrm{LE_{women}}-10}{\mathrm{LE_{women}}}\right),
\end{align*}
with $\mathrm{LE_{men}} = 14.45$ and $\mathrm{LE_{women}} = 22.55$ the life expectancies at the two retirement ages. Since $\epsilon$ depends on $\beta$, it is refreshed with $\beta$ inside the calibration. A variant sets the pre-reform $\epsilon$ to zero. The reform around $t=2010$ that raised $\epsilon$ is an experiment on the calibrated model (\cref{\refnumsec:shock}), not a calibration target.


\smalltitle{Calibration.}
Several of the model parameters are set directly using the data; what follows covers those that have to be solved for.

First, we calibrate the income distribution and relative working hours across the income distribution using $\eta_j$ and $X_j$. Specifically, let $z_j^{\eta}$ and $z_j^{x}$ denote income and working hours of type $j$ relative to the average \emph{formal} household. The reference is the same for all types, informal ones included, and the informal type is compared to it without being part of it, so that
\begin{align}\label{\refeq:calibration:z}
    \sum_{i}\gamma_i z_i^{\eta} = \sum_{i}\gamma_i z_i^{x} = 1,
\end{align}
where, as throughout, the sums over $i$ run over the formal types only. In the data this means dividing each group's income and hours by the population-weighted mean over the formal groups.

With $\Lambda_t\equiv w_t(1-\tau_t)+\theta_{t+1}\overline{b}_{t+1}/(R_{t+1}/p_t)$, formal labour supply \eqref{\refeq:formalOpt} is $h_{t,i} = (\eta_i/X_i)^{\xi}\Lambda_t^{\xi}$, so that $h_t = \Gamma_h\Lambda_t^{\xi}$ and $h_{t,i}/h_t = (\eta_i/X_i)^{\xi}/\Gamma_h$. The relative income and hours of the formal types are then
\begin{align*}
        z_i^{\eta} &= \dfrac{w_t(1-\tau_t)h_{t,i} \eta_i}{w_t(1-\tau_t)h_t} = \dfrac{\eta_i^{1+\xi}/X_i^{\xi}}{\sum_{j} \gamma_ j\eta_j^{1+\xi}/X_j^{\xi}},  \\
        z_i^{x} &= \frac{h_{t,i}}{\sum_j\gamma_j h_{t,j}} = \frac{(\eta_i/X_i)^{\xi}}{\sum_j \gamma_j(\eta_j/X_j)^{\xi}}.
\end{align*}
Both conditions are ratios and leave two normalizations free. The first is the level of the labour aggregator $\Gamma_h$. The second is the unit in which hours are counted: the rescaling $\eta_i\rightarrow\eta_i/\mu$, $X_i\rightarrow\mu^{-(1+\xi)/\xi}X_i$ leaves $\eta_i^{1+\xi}/X_i^{\xi}$ and therefore efficiency units $\eta_ih_{t,i}$, the aggregate $h_t$ and the equilibrium unchanged, and scales every $h_{t,i}$ by $\mu$. We spend both normalizations in the calibration stage:
\begin{align}\label{\refeq:calibration:yNorm}
    \Gamma_h\equiv\sum_i\gamma_i\frac{\eta_i^{1+\xi}}{X_i^{\xi}} = 1, \qquad \sum_i\gamma_i\left(\frac{\eta_i}{X_i}\right)^{\xi} = 1.
\end{align}
The second one is a real restriction on the informal calibration below, not a convenience. The aggregate $h_t$ weights hours by productivity, and average formal hours are
\begin{align}\label{\refeq:calibration:M}
    \sum_i\gamma_ih_{t,i} = M_t\,h_t, \qquad M_t \equiv \frac{\sum_i\gamma_i(\eta_i/X_i)^{\xi}}{\Gamma_h},
\end{align}
so under \eqref{\refeq:calibration:yNorm}, and only under it, $M_t = 1$ and $h_t$ equals average formal hours in the same clock unit as the individual $h_{t,j}$. The informal hours target below is stated against average formal hours, and carries $M_t$ so that it stays correct when the second normalisation is spent otherwise (the common-$X$ variant at the end of this section).

Using $y_i^{\eta} = \eta_i^{1+\xi}/X_i^{\xi}$ and $y_i^x = (\eta_i/X_i)^{\xi}$ as shorthand, we can set up the conditions as linear systems instead:
\begin{align*}
    \left(\bm{z}^{\eta} \bm{\gamma}'-\bm{I}\right)\bm{y}^{\eta} &= \bm{0} \\
    \left(\bm{z}^{x} \bm{\gamma}'-\bm{I}\right) \bm{y}^{x} &= \bm{0},
\end{align*}
where the bold notation indicates vectors and matrices over the formal types. The $\bm{y}$ vectors are eigenvectors with eigenvalue $\bm{\gamma}'\bm{z} = 1$. Since $\bm{z}\bm{\gamma}'$ has rank one, the eigenvector is $\bm{z}$ itself, and \eqref{\refeq:calibration:yNorm} fixes its scale at $\bm{\gamma}'\bm{y} = 1$, so $\bm{y}^{\eta} = \bm{z}^{\eta}$ and $\bm{y}^{x} = \bm{z}^{x}$. The vectors of productivity and disutility of labour follow as
\begin{align}
    \bm{\eta} &= \frac{\bm{y}^{\eta}}{\bm{y}^x} \label{\refeq:calibration:etai} \\
    \bm{X} &= \frac{\bm{\eta}}{(\bm{y}^x)^{1/\xi}}, \label{\refeq:calibration:Xi}
\end{align}
with division and powers taken element by element.

To capture the relative income and working hours of the informal households, requires a bit more. By \eqref{\refeq:calibration:z} and \eqref{\refeq:calibration:yNorm}, the ratios of income and working hours are
\begin{align*}
    z_0^{\eta} &= \frac{w_t^0h_{t,0}\eta_{0}}{w_t(1-\tau_t)h_t}, \\
    z_0^x &= \dfrac{h_{t,0}}{\sum_i \gamma_{i} h_{t,i}} = \frac{h_{t,0}}{M_th_t}.
\end{align*}
Write $k_t\equiv s_{t-1}/\nu_t$. In equilibrium $h_t = \Theta_{h,t}k_t^{\alpha\xi/(1+\alpha\xi)}$, so the factor prices \eqref{\refeq:factorPrices} and \eqref{\refeq:factorPrices0} give $w_t = (1-\alpha)\Theta_{h,t}^{-\alpha}k_t^{\alpha/(1+\alpha\xi)}$ and $w_t^0 = \left((1-\alpha)/\Gamma_{h,t}^{\alpha}\right)^{1/(1+\alpha\xi)}k_t^{\alpha/(1+\alpha\xi)}$. The ratio $z_0^{\eta}/z_0^x = M_t\eta_0w_t^0/\left(w_t(1-\tau_t)\right)$ then identifies $\eta_0$, and the informal labour supply condition \eqref{\refeq:informalOpt}, $h_{t,0} = (\eta_0w_t^0/X_0)^{\xi} = z_0^xM_t\Theta_{h,t}k_t^{\alpha\xi/(1+\alpha\xi)}$, identifies $X_0$. Here $k_t$ drops out of both:
\begin{align}
    \eta_{0} &= \frac{z_0^{\eta}}{z_0^x}\frac{(1-\alpha)(1-\tau_t)}{M_t\,\Theta_{h,t}^{\alpha}\left(\frac{1-\alpha}{\Gamma_{h,t}^{\alpha}}\right)^{\frac{1}{1+\alpha\xi}}} \label{\refeq:calibration:eta0}\\
    X_{0} &= \eta_{0}\frac{\left(\frac{1-\alpha}{\Gamma_{h,t}^{\alpha}}\right)^{\frac{1}{1+\alpha\xi}}}{\left(\Theta_{h,t}z_0^xM_t\right)^{1/\xi}} \label{\refeq:calibration:X0}
\end{align}
with $M_t = 1$ under \eqref{\refeq:calibration:yNorm}. Note that the product $\eta_0h_{t,0}$, and with it $\eta_0^{1+\xi}/X_0^{\xi}$ and everything the informal household contributes to the equilibrium, does not depend on $M_t$: the hours unit changes how the informal hours are counted, not what they earn.
We note that all of this is evaluated at $t=2010$, our baseline year. However, while we can identify most of the terms, the coefficient function $\Theta_{h,t}$ still has to be solved for us to get this. Thus, we add these two conditions determining $\eta_0, X_0$ as targets that we have to solve alongside the politico-economic equilibrium (to get $\Theta_{h,t}$ in the baseline year).



Given vectors of productivities and disutility of labour for formal households, $\bm{\eta}_i, \bm{X}_i$, we can determine $\theta$ from the ratio of replacement rates. In particular, we solve for $\theta$ (constant across $t$ in this calibration) from the requirement:
 \begin{align*}
        RR \equiv \dfrac{\theta+(1-\theta)\frac{h_t}{\eta_3 h_{t,3}}}{\theta+(1-\theta)\frac{h_t}{\eta_1 h_{t,1}}} = \frac{\theta + (1-\theta)\frac{\Gamma_{h,t} X_3^{\xi}}{\eta_3^{1+\xi}}}{\theta+(1-\theta)\frac{\Gamma_{h,t} X_1^{\xi}}{\eta_1^{1+\xi}}},
\end{align*}
where we have data on $RR$ from OECD data, and the two groups $j=1, 3$ represent half-mean and mean wages in our data for Argentina. Given that we know $\eta_i, X_i$ for formal households already, we can simply compute $\theta$ from
\begin{align}
    \theta = \frac{RR \frac{\Gamma_{h,t} X_1^{\xi}}{\eta_1^{1+\xi}} - \frac{\Gamma_{h,t} X_3^{\xi}}{\eta_3^{1+\xi}}}{1-\frac{\Gamma_{h,t} X_3^{\xi}}{\eta_3^{1+\xi}}-RR\left(1-\frac{\Gamma_{h,t} X_1^{\xi}}{\eta_1^{1+\xi}}\right)}
\end{align}



Finally, we need to determine parameters $\omega$ and discount parameters $\beta$. Note that we assume that they are identical across time and household types in the calibration. They are identified by ensuring that the politico-economic equilibrium replicates the tax rate $\tau_{t_0}$ derived above and a capital--output ratio of 3.23 in the baseline year 2010. Capital at $t$ is the previous generation's savings, $K_t = s_{t-1}/\nu_t$, so with Cobb--Douglas production the ratio is
\begin{align*}
    \frac{K_t}{Y_t} = \frac{s_{t-1}/\nu_t}{\left(s_{t-1}/\nu_t\right)^{\alpha} h_t^{1-\alpha}} = \frac{\alpha}{R_t},
\end{align*}
which is a ratio of a stock to \emph{one period's} output. Multiplying by the thirty years in a period puts it in the units the data are reported in. The savings rate
\begin{align*}
    \text{Savings rate}_t = \frac{s_t}{\left(s_{t-1}/\nu_t\right)^{\alpha} h_t^{1-\alpha}}
\end{align*}
is reported alongside it, but is not a target: with capital fully depreciating between periods it equals $K_{t+1}/Y_t$, which is a stock over a period's flow rather than the annual saving rate the national accounts measure.


Most parameters can be calibrated initially, assuming values for parameters $\xi, \rho$. There are, however, four parameters that cannot be calibrated without knowing the politico-economic equilibrium path: $\omega, \beta, \eta_0, X_0$. As the politico-economic equilibrium cannot be solved without these parameters, the final bit of calibration is solved as a nested fixed point problem. In \cref{\refnumsec} we describe this calibration procedure in more detail.

\smalltitle{The common-$X$ variant.} The calibration above uses relative hours as data and leaves the level of hours unidentified, which is the variant A of the rich-economy calibration in the US documentation. The variant B there is available here too, and is the variant the paper prints. Under it the formal types share one leisure parameter, $X_i = X$, and \eqref{\refeq:calibration:etai} collapses to the closed form
\begin{align}\label{\refeq:calibration:etaCommonX}
    \eta_i = \left(z_i^{\eta}\right)^{\frac{1}{1+\xi}}X^{\frac{\xi}{1+\xi}},
\end{align}
which imposes $\Gamma_h = 1$ for every $X$. The second normalisation in \eqref{\refeq:calibration:yNorm} is then not imposed; in its place the observed average workweek of formal workers, converted to a share of discretionary time, is imposed on average formal hours, $\sum_i\gamma_ih_{t_0,i} = \overline{h}_{t_0}$. Because $X$ enters no aggregate ($\eta_i^{1+\xi}/X_i^{\xi} = z_i^{\eta}$ whatever $X$ is), the nested fixed point is solved at $X = 1$ and $X$ follows in closed form afterwards, exactly as in the US documentation: with $M_{t_0} = X^{-\xi/(1+\xi)}\sum_i\gamma_i(z_i^{\eta})^{\xi/(1+\xi)}$ from \eqref{\refeq:calibration:M}, the target gives
\begin{align}\label{\refeq:calibration:Xsolve}
    X = \left(\frac{h_{t_0}}{\overline{h}_{t_0}}\sum_i\gamma_i\left(z_i^{\eta}\right)^{\frac{\xi}{1+\xi}}\right)^{\frac{1+\xi}{\xi}}.
\end{align}
The politico-economic equilibrium, $\beta$, $\omega$, the tax rate and the capital--output ratio are unchanged by the choice of variant; the informal parameters rescale as $\eta_0\rightarrow\eta_0/M_{t_0}$ and $X_0\rightarrow X_0M_{t_0}^{-1-1/\xi}$ by \eqref{\refeq:calibration:eta0}--\eqref{\refeq:calibration:X0}, and relative formal hours become a prediction, $h_{t,i}/\sum_j\gamma_jh_{t,j} \propto (z_i^{\eta})^{\xi/(1+\xi)}$, to be reported against the data $z_i^x$.
